\documentclass[../../main.tex]{subfiles}
\begin{document}
\section{子丛}\label{section:ch10-子丛}

设 $\pi\colon E\to M$ 为光滑 $k$ 秩向量丛。$E$ 的一个\textbf{光滑子丛} $D\subseteq E$ 是一个嵌入子流形（带边或不带边），在限制映射下它本身是 $M$ 上的光滑向量丛，且在每个 $p\in M$ 处的纤维 $D_p\subseteq E_p$ 是 $E_p$ 的一个 $m$ 维线性子空间。（这里 $m\leqslant k$ 是一个固定整数。）为了使限制映射成为向量丛投影，我们只要求 $D_p$ 是 $E_p$ 的线性子空间，并要求 $\pi|_D$ 是光滑向量丛投影；局部平凡化从 $E$ 的局部平凡化继承而来，只要它在每个子空间 $D_p$ 上是线性的。

\textbf{例 10.30}。若 $S\subseteq M$ 是带边或不带边的光滑嵌入子流形，则 $TS$ 自然是 $TM|_S$ 的一个光滑子丛。$\blacktriangleright$

\begin{figure}[htbp]
\centering
\includegraphics[width=0.85\textwidth]{fig-10-4.png}
\caption{向量丛的一个子丛}\label{figure:ch10-子丛}
\end{figure}

\begin{exercise}\label{exercise:ch10-31}
给定光滑向量丛 $E\to M$ 与光滑子丛 $D\subseteq E$，证明包含映射 $\iota\colon D\hookrightarrow E$ 是 $M$ 上的光滑丛同态。
\end{exercise}

下面的引理给出了一个方便的条件，用于检验线性子空间族 $\{D_p\subseteq E_p:p\in M\}$ 是否构成光滑子丛。

\begin{lemma}[子丛的局部标架判别法]\label{lemma:ch10-子丛局部标架判别}
设 $\pi\colon E\to M$ 为光滑向量丛，并假设对每个 $p\in M$ 给定一个 $m$ 维线性子空间 $D_p\subseteq E_p$。则 $D=\bigcup_{p\in M}D_p\subseteq E$ 是 $E$ 的光滑子丛，当且仅当下面的条件满足：
\begin{align}
&\text{$M$ 的每一点都有一个邻域 $U$，在其上存在光滑局部截面 }\sigma_1,\dots,\sigma_m\colon U\to E,\nonumber\\
&\text{使得对每个 }q\in U,\ \sigma_1(q),\dots,\sigma_m(q)\text{ 构成 }D_q\text{ 的一组基。}\label{align:ch10-子丛判别}
\end{align}
\end{lemma}

\begin{proof}
若 $D$ 是光滑子丛，则按定义每个 $p\in M$ 都有一个邻域 $U$，在其上存在 $D$ 的光滑局部平凡化，而\cref{example:ch10-局部标架}~表明在每个这样的 $U$ 上存在 $D$ 的光滑局部标架。这种局部标架按定义就是一组光滑截面 $\tau_1,\dots,\tau_m\colon U\to D$，它们在每点 $p\in U$ 的像构成 $D_p$ 的一组基。我们所求的 $E$ 的光滑截面由与包含映射 $\iota\colon D\hookrightarrow E$ 复合得到：$\sigma_j=\iota\circ\tau_j$。

反过来，假设 $E\to M$ 是秩 $k$ 的光滑丛，且 $D\subseteq E$ 满足\eqref{align:ch10-子丛判别}。选取 $p$ 的一个邻域 $U$ 以及 $E$ 在 $U$ 上的一个光滑局部平凡化 $\Phi\colon\pi^{-1}(U)\to U\times\mathbb{R}^k$。条件\eqref{align:ch10-子丛判别}给出 $D$ 在 $U$ 上的光滑截面 $\sigma_1,\dots,\sigma_m$。在 $\Phi$ 下，这些截面成为 $U\times\mathbb{R}^k$ 的光滑截面，我们可在 $U$ 上把它们补成 $\mathbb{R}^k$ 的一组光滑局部标架 $(\widehat\sigma_1,\dots,\widehat\sigma_k)$，使得对每个 $q\in U$，$\widehat\sigma_1(q),\dots,\widehat\sigma_m(q)$ 恰为 $\sigma_1(q),\dots,\sigma_m(q)$。在以这个标架为坐标的新局部平凡化下，$\Phi$ 把 $D\cap\pi^{-1}(U)$ 映到子集
\begin{align*}
\big\{(q,(s_1,\dots,s_m,0,\dots,0)):q\in U\big\}\subseteq U\times\mathbb{R}^k,
\end{align*}
这是一个嵌入子流形（若 $U$ 带边则亦带边）。此外，由
\begin{align*}
\Psi\left(\sum_{i=1}^m s_i\,\sigma_i(q)\right)=(q,(s_1,\dots,s_m))
\end{align*}
定义的映射 $\Psi\colon D\cap\pi^{-1}(U)\to U\times\mathbb{R}^m$ 是 $D$ 的一个光滑局部平凡化，故 $D$ 本身是一个光滑向量丛。
\end{proof}

\begin{example}[子丛]\label{example:ch10-子丛例子}
(a) 若 $M$ 是光滑流形，而 $V$ 是 $M$ 上处处非零的光滑向量场，则集合 $D\subseteq TM$（其在每点 $p\in M$ 处的纤维为 $V_p$ 的线性张成）是 $TM$ 的一个光滑 $1$ 维子丛。

(b) 假设 $E\to M$ 是任意平凡丛，而 $(E_1,\dots,E_k)$ 是 $E$ 的一个光滑整体标架。若 $0\leqslant m\leqslant k$，则由 $D_p=\mathrm{span}(E_1|_p,\dots,E_m|_p)$ 对每个 $p\in M$ 定义的子集 $D\subseteq E$ 是 $E$ 的一个光滑子丛。

(c) 假设 $M$ 是带边或不带边的光滑流形，$S\subseteq M$ 是带边或不带边的浸入 $k$ 子流形。习题~10-14 要求你证明 $TS$ 是环境切丛 $TM|_S$ 的一个光滑 $k$ 秩子丛。
\end{example}

下面的定理表明如何得到更多的子丛。设 $E\to M$ 与 $E'\to M$ 是向量丛，而 $F\colon E\to E'$ 是 $M$ 上的丛同态。对每个 $p\in M$，线性映射 $F|_{E_p}$ 的秩称为 $F$ 在 $p$ 处的秩。若 $F$ 在所有 $p\in M$ 处的秩都相同，则称 $F$ 具有\textbf{常秩}。

\begin{theorem}[常秩丛同态的核与像]\label{theorem:ch10-常秩丛同态}
设 $E$ 与 $E'$ 是光滑流形 $M$ 上的光滑向量丛，$F\colon E\to E'$ 是 $M$ 上的光滑丛同态。定义子集 $\mathrm{Ker}\,F\subseteq E$ 与 $\mathrm{Im}\,F\subseteq E'$ 为
\begin{align*}
\mathrm{Ker}\,F=\bigcup_{p\in M}\mathrm{Ker}\big(F|_{E_p}\big),\qquad
\mathrm{Im}\,F=\bigcup_{p\in M}\mathrm{Im}\big(F|_{E_p}\big).
\end{align*}
则 $\mathrm{Ker}\,F$ 与 $\mathrm{Im}\,F$ 分别是 $E$ 与 $E'$ 的光滑子丛，当且仅当 $F$ 具有常秩。
\end{theorem}

\begin{proof}
一个方向是显然的：由于丛的纤维维数处处相同，$\mathrm{Ker}\,F$ 与 $\mathrm{Im}\,F$ 成为子丛当然需要常秩条件。为证充分性，设 $F$ 有常秩 $r$，并令 $k$ 与 $k'$ 分别为丛 $E$ 与 $E'$ 的秩。任取 $p\in M$，在 $p$ 的某邻域 $U$ 上选取 $E$ 的一个光滑局部标架 $(\sigma_1,\dots,\sigma_k)$。对每个 $i$，映射 $F\circ\sigma_i\colon U\to E'$ 是 $E'$ 的光滑局部截面，这些截面张成 $(\mathrm{Im}\,F)|_U$。必要时重排指标，可假设 $\{F\circ\sigma_1(p),\dots,F\circ\sigma_r(p)\}$ 构成 $\mathrm{Im}(F|_{E_p})$ 的一组基，由连续性它们在 $p$ 的某邻域 $U'$ 上仍线性无关。由于 $F$ 有常秩，这意味着 $(F\circ\sigma_1,\dots,F\circ\sigma_r)$ 在 $U'$ 上构成 $\mathrm{Im}\,F$ 的光滑局部标架。由于在每点附近都可如此做，局部标架判别法表明 $\mathrm{Im}\,F$ 是 $E'$ 的光滑子丛。

为证 $\mathrm{Ker}\,F$ 也是光滑子丛，令 $U'$ 与 $(\sigma_i)$ 如上，并令 $V\subseteq E|_{U'}$ 为由 $\sigma_1,\dots,\sigma_r$ 张成的光滑子丛。光滑丛同态 $F|_V\colon V\to(\mathrm{Im}\,F)|_{U'}$ 是双射，因而由\cref{proposition:ch10-光滑丛同态性质}~是光滑丛同构。定义光滑丛同态 $P\colon E|_{U'}\to E|_{U'}$ 为 $P(v)=v-(F|_V)^{-1}\circ F(v)$。若 $v\in V$，则 $F(v)=(F|_V)(v)$，于是
\begin{align*}
F\big(P(v)\big)=F(v)-F\circ(F|_V)^{-1}\circ(F|_V)(v)=0.
\end{align*}
另一方面，若 $v\in\mathrm{Ker}\,F$，则 $P(v)=v$，故仍有 $F(P(v))=F(v)=0$。由于 $V$ 与 $(\mathrm{Ker}\,F)|_{U'}$ 一起张成 $E|_{U'}$，可知 $P$ 取值于 $(\mathrm{Ker}\,F)|_{U'}$，又因它在 $(\mathrm{Ker}\,F)|_{U'}$ 上限制为恒等映射，其像恰为 $(\mathrm{Ker}\,F)|_{U'}$。于是 $P$ 有常秩，由上一段的论证，$(\mathrm{Ker}\,F)|_{U'}=\mathrm{Im}\,P$ 是 $E|_{U'}$ 的光滑子丛。由于在每点附近都可如此做，$\mathrm{Ker}\,F$ 是 $E$ 的光滑子丛。
\end{proof}

下面的引理给出了在 $\mathbb{R}^n$ 的子流形上构造切丛的更多光滑子丛的另一种方法。

\begin{lemma}[正交补丛]\label{lemma:ch10-正交补丛}
设 $M$ 是 $\mathbb{R}^n$ 中带边或不带边的浸入子流形，$D$ 是 $TM|_M$（即 $T\mathbb{R}^n|_M$）的一个光滑 $k$ 秩子丛。对每个 $p\in M$，令 $D_p^\perp$ 表示 $D_p$ 在 $T_p\mathbb{R}^n$ 中关于 Euclid 点积的正交补，并令 $D^\perp\subseteq T\mathbb{R}^n|_M$ 为子集
\begin{align*}
D^\perp=\{(p,v)\in T\mathbb{R}^n:p\in M,\ v\in D_p^\perp\}.
\end{align*}
则 $D^\perp$ 是 $T\mathbb{R}^n|_M$ 的一个光滑 $(n-k)$ 秩子丛。对每个 $p\in M$，在 $p$ 的某邻域上存在 $D^\perp$ 的光滑正交标架。
\end{lemma}

\begin{proof}
任取 $p\in M$，并令 $(X_1,\dots,X_k)$ 为 $D$ 在 $M$ 中 $p$ 的某邻域 $V$ 上的光滑局部标架。由于浸入子流形是局部嵌入的，必要时缩小 $V$，可设它是某个坐标球或半球 $U\subseteq\mathbb{R}^n$ 中的一个切片。由于 $V$ 在 $U$ 中闭，\cref{proposition:ch08-11}~表明我们可把 $(X_1,\dots,X_k)$ 扩充为 $T\mathbb{R}^n$ 在 $U$ 上的光滑局部标架 $(\widetilde X_1,\dots,\widetilde X_n)$，再由\cref{lemma:ch08-13}~得到 $U$ 上的光滑正交标架 $(E_j)$，使得对每个 $p\in U$ 有 $\mathrm{span}(E_1|_p,\dots,E_k|_p)=\mathrm{span}(X_1|_p,\dots,X_k|_p)=D_p$。于是 $(E_{k+1},\dots,E_n)$ 限制为 $D^\perp$ 在 $V$ 上的光滑正交标架。因此 $D^\perp$ 满足局部标架判别法，从而是 $T\mathbb{R}^n|_M$ 的光滑子丛。
\end{proof}

\begin{corollary}[$\mathbb{R}^n$ 子流形的法丛]\label{corollary:ch10-法丛}
若 $M\subseteq\mathbb{R}^n$ 是带边或不带边的浸入 $m$ 维子流形，则其法丛 $NM$ 是 $T\mathbb{R}^n|_M$ 的一个光滑 $(n-m)$ 秩子丛。对每个 $p\in M$，在 $p$ 的某邻域上存在 $NM$ 的光滑正交标架。
\end{corollary}

\begin{proof}
把引理~\ref{lemma:ch10-正交补丛}~应用于光滑子丛 $TM\subseteq T\mathbb{R}^n|_M$。
\end{proof}
\end{document}
